\section{Asymptotic quantities and the counterexamples}
\label{sec:counterexamples}

Fix the cutoff of \Cref{sec:construction}. Throughout this section \(L\)
is above a single threshold for the preceding construction, uniformly for
\(s\in[0,1]\). Write \(H\) and \(N\) for the source and Newton potential
at the fixed point. All constants below are independent of these two
parameters once that threshold is fixed.

\subsection{Mass, momentum, and charges}

\begin{proposition}\label{prop:charges}
The constructed data have finite asymptotic quantities
\begin{equation}\label{eq:vanishing-charges}
 P_{\mathrm{ADM}}=0,\qquad Q_e=Q_b=0,
\end{equation}
and
\begin{equation}\label{eq:mass}
 m=2+2N(0)+\frac1{2\pi}\int_{\R^3}H(x)\,\dd x
   =2+O(s^2/L).
\end{equation}
In particular \(m\ge2\). If \(M\) is a uniform upper bound for \(u\), then
\begin{equation}\label{eq:mass-lower-bound}
 m-2\ge\frac{s^2}{3M^3L}.
\end{equation}
The angular momentum in \eqref{eq:adm-angular} is
\begin{equation}\label{eq:angular-momentum}
 J=-\frac{4s^2}{15}.
\end{equation}
For \(s>0\), neither electromagnetic field is identically zero.
\end{proposition}

\begin{proof}
The Cartesian decay \(K=O(r^{-3})\), with \(\tau=0\), makes the integrals
in \eqref{eq:adm-momentum} of order \(O(r^{-1})\). This proves the first
part of \eqref{eq:vanishing-charges}.

The conformal weights leave both Maxwell flux forms unchanged. On every
coordinate sphere, the electric flux is
\[
 \int_{S_r}\iota_{\mathcal E}\dd V_g
 =\int_{S_r}\dd(f\,\dd\phi)
 =2\pi\bigl(f(r,\pi)-f(r,0)\bigr)=0.
\]
The identical calculation for \(h\) gives zero magnetic flux. Equivalently,
these are integrals of globally smooth exact two-forms over closed
surfaces. The charge limits are therefore zero.

Substituting \(g_{ij}=u^4\delta_{ij}\) in \eqref{eq:adm-energy} gives
\begin{equation}\label{eq:conformal-mass-flux}
 E_{\mathrm{ADM}}
 =-\frac1{2\pi}\lim_{R\to\infty}
   \int_{S_R}u^3u_r\,\dd A_\delta.
\end{equation}
Since \(u-1=O(R^{-1})\) and \(u_r=O(R^{-2})\), replacing \(u^3\) by one
has an integrated error \(O(R^{-1})\). The source \(H\) is integrable and
the ordinary Newton potential satisfies the exact flux identity
\[
 \int_{S_R}N_r\,\dd A_\delta=-\int_{B_R}H\,\dd x.
\]
For the reflected term, smoothness of \(N\) at the origin gives
\[
 R^{-1}N(\omega/R)=\frac{N(0)}{R}+O(R^{-2}),
 \qquad
 \partial_R\bigl(R^{-1}N(\omega/R)\bigr)
       =-\frac{N(0)}{R^2}+O(R^{-3}).
\]
Together with \(u_0=1+1/r\), these identities prove the existence of the
limit and the exact formula \eqref{eq:mass}. In particular no first-moment
assumption on \(H\) or multipole expansion of the ordinary potential is
needed. The estimates of \Cref{prop:conformal-factor} give
\(N(0)=O(s^2/L^2)\), \(\int H=O(s^2/L)\), and all contributions in
\eqref{eq:mass} are nonnegative. Since the linear momentum is zero,
\(m=E_{\mathrm{ADM}}\).

On \(r\ge2L\) the cutoff is one, and
\begin{equation}\label{eq:radial-tails}
 \widetilde E=\frac{2s\cos\theta}{r^2}\partial_r,
 \qquad
 \widetilde B=\frac{s(3\cos^2\theta-1)}{r^2}\partial_r.
\end{equation}
Using just the electric field and \(u\le M\),
\[
 H\ge\ell u^{-3}\ge\frac{s^2\cos^2\theta}{M^3r^4}
 \qquad(r\ge2L).
\]
As \(\int_{\Sph^2}\cos^2\theta\,\dd A=4\pi/3\), this gives
\[
 \int H\,\dd x\ge\frac{2\pi s^2}{3M^3L}.
\]
The mass formula proves \eqref{eq:mass-lower-bound}.

Finally \(\nu_r=u^{-2}\partial_r\), \(\dd A_g=u^4\dd A_\delta\), and
\(K=u^{-2}\sigma\), so all conformal weights in the angular-momentum
integrand cancel. The gravitational flux at every coordinate sphere is
\begin{align}
 J_{\mathrm{grav}}(r)
  &:=\frac1{8\pi}\int_{S_r}K(\nu_r,\eta)\,\dd A_g\notag\\
  &=\frac1{8\pi}\int_{S_r}\sigma(\partial_r,\partial_\phi)
           \,\dd A_\delta
   =-\frac{s^2c(r)^2}{4}\int_0^\pi\sin^5\theta\,\dd\theta
   =-\frac{4s^2c(r)^2}{15}.
   \label{eq:finite-radius-angular}
\end{align}
It is constant for \(r\ge2L\), proving \eqref{eq:angular-momentum} and
existence of the prescribed limit. The two nonzero angular polynomials in
\eqref{eq:radial-tails}, multiplied by the positive factor \(u^{-6}\), also
prove the final assertion.
\end{proof}

\subsection{The conserved electromagnetic correction}

The flux \eqref{eq:finite-radius-angular} vanishes near the horizon but
not at infinity. The following calculation identifies the boundary term
responsible for the difference, using precisely our constraint signs.

\begin{proposition}\label{prop:corrected-flux}
Let \(A_B=h\,\dd\phi\), the smooth magnetic potential one-form in
\Cref{sec:construction}. Then
\begin{equation}\label{eq:corrected-divergence}
 \operatorname{div}_g\left(K(\cdot,\eta)^{\sharp_g}
                              +2h\mathcal E\right)=0.
\end{equation}
The electromagnetic boundary contribution on every coordinate sphere is
\begin{equation}\label{eq:electromagnetic-correction}
 J_{\mathrm{EM}}(r)
 :=\frac1{4\pi}\int_{S_r}h\,\iota_{\mathcal E}\dd V_g
 =\frac{4s^2c(r)^2}{15}.
\end{equation}
Consequently \(\mathcal J(r):=J_{\mathrm{grav}}(r)+J_{\mathrm{EM}}(r)=0\)
for all \(r\ge1\), whereas \(J_{\mathrm{EM}}(r)\to4s^2/15\) at infinity.
\end{proposition}

\begin{proof}
Since \(A_B\) is invariant and \(A_B(\eta)=h\), Cartan's identity gives
\[
 \iota_\eta\dd A_B=-\dd h.
\]
Using \(\dd A_B=\iota_{\mathcal B}\dd V_g\) and evaluating on
\(\mathcal E\) gives
\[
 \mathcal E(h)=-\dd V_g(\mathcal E,\mathcal B,\eta).
\]
The symmetry of \(K\) and the Killing equation imply
\[
 \operatorname{div}_g\bigl(K(\cdot,\eta)^{\sharp_g}\bigr)
   =(\operatorname{div}_gK)(\eta)
   =2\dd V_g(\mathcal E,\mathcal B,\eta).
\]
Combining this with \(\operatorname{div}_g\mathcal E=0\) proves
\eqref{eq:corrected-divergence}. Directly from the flux form,
\[
 J_{\mathrm{EM}}(r)
  =\frac12\int_0^\pi h f_\theta\,\dd\theta
  =s^2c(r)^2\int_0^\pi\sin^3\theta\cos^2\theta\,\dd\theta
  =\frac{4s^2c(r)^2}{15}.
\]
Now apply \eqref{eq:finite-radius-angular}.
\end{proof}

There is no singular gauge choice in this computation: \(A_B\) is a
globally smooth one-form. The scalar \(h=A_B(\eta)\) vanishes on the
axis, where \(\eta=0\). An additive constant \(C\) in a scalar potential
would change the boundary contribution only by \(C Q_e=0\).
The norm of \(A_B\) is \(O(r^{-1})\), while
\(A_B(\eta)=h\) need not tend to zero because \(\abs\eta=O(r)\).
Moreover, for \(s>0\), the zero-charge fields in \eqref{eq:radial-tails} retain
angularly varying \(r^{-2}\) terms. Multiplication by
\(u^{-6}=1+O(r^{-1})\) preserves these leading terms. They satisfy
\eqref{eq:falloff}, but they are not \(O(r^{-3})\).
The divergence constraint does not prohibit this: in Euclidean space,
\[
 \operatorname{div}_\delta\bigl(r^{-2}a(\omega)\partial_r\bigr)=0
\]
for every smooth angular function \(a\). No curl-free or stationarity
condition is imposed. Thus \Cref{prop:corrected-flux} is consistent with
electromagnetic angular-momentum conservation; it is the identification
of that conserved flux with \eqref{eq:adm-angular} which fails here.

\subsection{Strict violation and nondegenerate equality}

The area estimate and angular momentum already give, uniformly in the
family,
\begin{equation}\label{eq:strict-branch}
 A\ge64\pi>\frac{32\pi}{15}\ge8\pi\abs J.
\end{equation}
Since \(Q=0\), this is the strict physical area branch. Set
\begin{equation}\label{eq:defect}
 D(s,L)=m(s,L)^2-\frac{A(s,L)}{16\pi}
                         -\frac{4\pi J(s,L)^2}{A(s,L)}.
\end{equation}
By \eqref{eq:area-estimate}, \eqref{eq:mass}, and
\eqref{eq:angular-momentum},
\begin{equation}\label{eq:negative-limit}
 \lim_{L\to\infty}D(1,L)=-\frac1{225}.
\end{equation}
Hence \(D(1,L)<0\) for all sufficiently large finite \(L\).

\begin{proposition}\label{prop:equality}
For every sufficiently large fixed \(L\), there exists \(s_*\in(0,1)\)
such that \(D(s_*,L)=0\). These equality data satisfy the strict branch
\eqref{eq:strict-branch} and have no Kerr--Newman realization with the
specified asymptotic parameters and induced fields.
\end{proposition}

\begin{proof}
Write the uniform area estimate as
\[
 64\pi\le A\le64\pi+C_A s^2/L^2.
\]
The lower mass bound \eqref{eq:mass-lower-bound} and \(A\ge64\pi\) imply
\begin{equation}\label{eq:positive-small-amplitude}
 D(s,L)\ge
 \frac{4s^2}{3M^3L}-\frac{C_A s^2}{16\pi L^2}-\frac{s^4}{225}.
\end{equation}
Choose \(L\) above the thresholds for the previous propositions,
\eqref{eq:negative-limit}, and positivity of the quadratic coefficient
on the right of \eqref{eq:positive-small-amplitude}. For this fixed \(L\),
the right side is positive at sufficiently small nonzero \(s\).

The fixed point depends continuously on \(s\) in the supremum norm.
Restriction to \(S\) therefore gives continuity of \(A\), and the
integrable source bounds give continuity of \(N(0)\) and \(\int H\).
The mass formula and \eqref{eq:angular-momentum} show that \(D(s,L)\) is
continuous. Its positive value at a nonzero small parameter and its
negative value at \(s=1\) yield a root \(s_*\in(0,1)\) by the
intermediate value theorem.

At this root, \(Q_e=Q_b=0\), \(J\ne0\), and both electromagnetic fields
are nonzero. The strict branch and the equality equation give
\begin{equation}\label{eq:subextremality}
 m^2-\abs J=\frac{(A-8\pi\abs J)^2}{16\pi A}>0.
\end{equation}
Thus the Kerr--Newman member with these parameters would be a
subextremal zero-charge member, namely Kerr with zero Maxwell tensor.
Indeed the dyonic Kerr--Newman potential is linear in its electric and
magnetic charges; see \cite[Equation~(3)]{CKST2018}. When both vanish,
\(F=0\). Every spacelike slice then has
\[
 (\iota_nF)|_{T\Omega}=0,\qquad
 \star_g(F|_{T\Omega})=0,
\]
independently of the slice or its future unit normal. This contradicts
the nonzero constructed fields, and proves the asserted obstruction.
\end{proof}

\begin{proof}[Proof of \Cref{thm:main}]
\Cref{lem:seeds,lem:conformal,prop:conformal-factor,prop:geometry}
provide smooth maximal electrovacuum data with every topological,
asymptotic, boundary, outermostness, and area condition in
\Cref{sec:conventions}. \Cref{prop:charges} supplies the finite ADM and
charge limits, positive mass and zero linear momentum. The physical branch
is strict by \eqref{eq:strict-branch}. For \(s=1\) and sufficiently large
finite \(L\), \eqref{eq:negative-limit} gives the first conclusion.
For the same large \(L\), \Cref{prop:equality} gives the second.
\end{proof}
